\section{Seven Lessons from Human Translation for Animal Decoding}
\label{sec:lessons}

The previous sections catalogued tools and audited their transfer. This
section distills the exchange into seven concrete lessons---what the history
of human machine translation, read as a computer-science narrative, teaches
anyone attempting to decode an animal. Each lesson pairs a human-MT finding
with a prescription for animal work. Table~\ref{tab:lessons} collects them.

\subsection{Lesson 1: Pretrain on unlabeled signal before anything else}
The single most robust finding of the last decade is that a large unlabeled
corpus, pretrained with a masked or contrastive objective, beats a clever
architecture trained from scratch~\cite{devlin2019bert,baevski2020wav2vec2,liu2019roberta}.
The animal prescription is immediate: invest first in a large, open,
unlabeled corpus of dog vocalizations and a self-supervised encoder, exactly
as AVES and NatureLM-audio have done for animals
generally~\cite{hagiwara2023aves,robinson2024naturelm}. Architecture is not the
bottleneck; data regime is.

\subsection{Lesson 2: Discretize before you translate}
BPE showed that composing rare words from frequent subword units unlocks
generalization~\cite{sennrich2016bpe}; neural codecs showed the same for
audio~\cite{zeghidour2021soundstream,defossez2022encodec}; and Project CETI's
``phonetic alphabet'' showed it for whales~\cite{sharma2024whale}. The
prescription: before claiming to translate, discover whether barks admit a
small set of reusable, recombinable units. If they do not, the signal is not
language-like and the translation framing should be dropped.

\subsection{Lesson 3: Bootstrap when no parallel corpus exists}
Human MT learned to manufacture its own supervision:
back-translation~\cite{sennrich2016improving,edunov2018backtrans} and fully
unsupervised translation~\cite{artetxe2018unsupervised,lample2018unsupervised}
removed the need for a bespoke parallel corpus. Animals have no parallel
corpus and never will. The prescription is to build a synthetic loop:
generate candidate vocalizations, play them back, and score the animal's
response---the generative-plus-reinforcement recipe of
Section~\ref{sec:transfer}. The error signal must come from the animal, not
from a human annotator's intuition.

\subsection{Lesson 4: Share structure across related systems}
Multilingual models showed that one encoder can serve hundreds of languages,
with low-resource pairs benefiting from high-resource
relatives~\cite{johnson2017multilingual,costajussa2022nllb,fan2021beyond}. The
prescription: train across species and across individuals, letting wolves,
coyotes, and many dogs inform a shared canid encoder, rather than fitting a
bespoke model per dog. Cross-species transfer is the animal analogue of
low-resource MT, and it is already paying off~\cite{ghani2023perch,babu2022xlsr}.

\subsection{Lesson 5: Evaluation precedes the claim}
Human MT earned its credibility through reference corpora and shared
tasks---BLEU, WMT, learned metrics validated against expert
judgment~\cite{papineni2002bleu,bojar2018wmt,rei2020comet,freitag2021experts}.
The animal field has no reference and therefore no comparable discipline. The
prescription is the demarcation of Section~\ref{sec:eval}: define the task,
disclose the protocol, and refuse the word ``translator'' until an independent
error signal exists. A field that cannot measure cannot improve.

\subsection{Lesson 6: Grounding is the wall, not scale}
Scaling laws make returns predictable~\cite{kaplan2020scaling,hoffmann2022chinchilla},
which tempts the inference that a big enough model will eventually translate a
bark. But the symbol-grounding problem is not a data deficit; it is a
referential mismatch~\cite{harnad1990grounding,bisk2020experience}. No amount
of compute creates a shared lexicon where none exists. The prescription:
spend effort on grounding and validation, not on parameter count, for the
upper rungs. Scale helps Rungs~1--3; it cannot fix a category error at
Rung~5.

\subsection{Lesson 7: Separate the interface from the measurement}
Modern systems render fluent sentences with a language
model~\cite{brown2020gpt3,ouyang2022instruct}, and fluency is easily mistaken
for fidelity---the hallucination and faithfulness literature documents exactly
this failure~\cite{ji2023hallucination,maynez2020faithfulness}. The
prescription: treat the rendered sentence as a user-interface layer, clearly
labeled as interpretation, and never let its fluency count as evidence. The
honest number lives upstream, at the classifier; the sentence is packaging.

\begin{table}[htbp]
\centering
\caption{Seven lessons from human machine translation, each paired with a
prescription for animal decoding.}
\label{tab:lessons}
\small
\setlength{\tabcolsep}{4pt}
\begin{tabular}{@{}p{0.6cm}p{4.7cm}p{8.0cm}@{}}
\toprule
\textbf{\#} & \textbf{Human-MT finding} & \textbf{Prescription for animals} \\
\midrule
1 & pretraining beats architecture~\cite{devlin2019bert,baevski2020wav2vec2} & build a large unlabeled bark corpus + SSL encoder first \\
2 & subwords / codecs unlock generalization~\cite{sennrich2016bpe,sharma2024whale} & discover reusable bark units before claiming translation \\
3 & unsupervised / back-translation~\cite{artetxe2018unsupervised,sennrich2016improving} & bootstrap supervision via playback + RL, not annotators \\
4 & multilingual sharing helps low-resource~\cite{costajussa2022nllb,johnson2017multilingual} & one shared canid encoder across species and individuals \\
5 & references + shared tasks = credibility~\cite{papineni2002bleu,bojar2018wmt} & no error signal, no ``translator''; disclose protocols \\
6 & scaling laws predict returns~\cite{kaplan2020scaling} & scale helps Rungs 1--3; grounding, not scale, blocks Rung 5 \\
7 & fluency $\neq$ fidelity~\cite{ji2023hallucination,maynez2020faithfulness} & label the sentence as interface; score the classifier \\
\bottomrule
\end{tabular}
\end{table}